% BotWork — Заявка «Студенческий Стартап» (VII очередь)
% Компиляция из каталога include/:
%   pdflatex -output-directory=../docs BotWork_Application.tex

\documentclass[aspectratio=169, 11pt]{beamer}

\usepackage[T2A]{fontenc}
\usepackage[utf8]{inputenc}
\usepackage[russian]{babel}
\usepackage{booktabs}

% Цвета
\definecolor{cBlue}{HTML}{B5D8E1}
\definecolor{cPink}{HTML}{FBB3C1}
\definecolor{cRed}{HTML}{FD493D}
\definecolor{cOrange}{HTML}{FCAF17}
\definecolor{cPurple}{HTML}{8F00FF}
\definecolor{cDark}{HTML}{1E1E1E}
\definecolor{cGray}{HTML}{555555}

% Тема
\usetheme{Madrid}
\usecolortheme{default}
\setbeamercolor{structure}{fg=cPurple}
\setbeamercolor{frametitle}{bg=cRed, fg=white}
\setbeamercolor{block title}{bg=cPurple, fg=white}
\setbeamercolor{block body}{bg=cBlue!15, fg=cDark}
\setbeamercolor{item}{fg=cPurple}
\setbeamerfont{frametitle}{series=\bfseries}
\setbeamertemplate{navigation symbols}{}
\setbeamertemplate{footline}[frame number]

% Декоративная полоса
\newcommand{\cbar}{%
  \textcolor{cPink}{\rule{0.2\linewidth}{2pt}}%
  \textcolor{cPurple}{\rule{0.2\linewidth}{2pt}}%
  \textcolor{cRed}{\rule{0.2\linewidth}{2pt}}%
  \textcolor{cOrange}{\rule{0.2\linewidth}{2pt}}%
  \textcolor{cBlue}{\rule{0.2\linewidth}{2pt}}\par\medskip}

% Руб
\newcommand{\R}{\,руб.}

\title{\textcolor{cPurple}{\Huge BotWork}}
\subtitle{Платформа для создания Telegram-ботов для бизнеса}
\author{Заявитель: \_\_\_\_\_\_\_\_\_\_\_\_\_\_}
\institute{Студенческий Стартап --- VII очередь}
\date{}

\begin{document}

% === СЛАЙД 1: Титульный ===
\begin{frame}[plain]
\titlepage
\centering\cbar
\end{frame}

% === СЛАЙД 2: О продукте (проблема + описание) ===
\begin{frame}{О продукте}
\cbar
\textbf{1. Какую проблему решает продукт}

\smallskip
Более \textbf{50--70 тыс.\ предпринимателей} в РФ ведут продажи через Telegram, но не имеют навыков для создания автоматизированных ботов. Платформы-конструкторы (BotHelp, SendPulse) стоят 10--70 тыс.\R\ и сложны.

\medskip
\textbf{2. Описание конечного продукта}

\smallskip
BotWork --- Telegram-боты \guillemotleft{}под ключ\guillemotright{}:
\begin{itemize}\setlength\itemsep{0pt}
  \item Автоматическая воронка: приветствие $\to$ презентация $\to$ оплата
  \item Интеграция платежей (ЮKassa, СБП, Crypto)
  \item Настройка под нишу (курсы, товары, услуги)
  \item Визуальный конструктор для проектирования логики бота
  \item \textcolor{cRed}{\textbf{Цена: 5--25 тыс.\R\ Срок: 3--7 дней. Поддержка 2 нед. бесплатно}}
\end{itemize}
\end{frame}

% === СЛАЙД 3: О продукте (область применения) ===
\begin{frame}{О продукте}
\cbar
\textbf{3. Область применения продукта}

\medskip
\textbf{\textcolor{cGray}{Способ применения:}} Бизнес получает готового бота, работающего 24/7 внутри Telegram.

\medskip
\textbf{\textcolor{cGray}{Ниши:}}
\begin{itemize}\setlength\itemsep{0pt}
  \item \textbf{Образование} --- курсы, вебинары, консультации через воронки
  \item \textbf{Услуги} --- запись клиентов (салоны, репетиторы), приём заявок
  \item \textbf{Товары} --- мини-магазины с каталогом, корзиной и оплатой
\end{itemize}

\medskip
\textbf{\textcolor{cGray}{Выгоды:}}
\begin{itemize}\setlength\itemsep{0pt}
  \item Автоматизация продаж, снижение нагрузки на персонал
  \item Увеличение конверсии на 15--25\%
  \item Всё внутри привычного Telegram --- без отдельных приложений
\end{itemize}
\end{frame}

% === СЛАЙД 4: Технологичность (техническое решение) ===
\begin{frame}{Технологичность стартап-проекта}
\cbar
\textbf{1. Техническое решение, инновационность}

\medskip
\textbf{\textcolor{cGray}{Стек:}}
\begin{itemize}\setlength\itemsep{0pt}
  \item Backend: Python + aiogram (асинхронный Telegram Bot API)
  \item Frontend-конструктор: HTML/CSS/JS, drag \& drop
  \item Модульная архитектура: блоки --- приветствие, кнопки, оплата, сбор данных, условия
  \item Интеграции: ЮKassa API, СБП, Crypto
\end{itemize}

\medskip
\textbf{\textcolor{cGray}{Инновационность:}}

Гибридная модель: клиент проектирует логику бота в визуальном конструкторе, затем передаёт разработчику для профессиональной реализации. Сочетание no-code с персональной разработкой \guillemotleft{}под ключ\guillemotright{} --- модель, которую не предлагают конкуренты.
\end{frame}

% === СЛАЙД 5: Технологичность (преимущества) ===
\begin{frame}{Технологичность стартап-проекта}
\cbar
\textbf{2. Преимущества технического решения}

\begin{itemize}\setlength\itemsep{3pt}
  \item \textbf{Конструктор снижает порог входа:} клиент описывает бота визуально, без технических знаний
  \item \textbf{Шаблонный подход:} типовые боты за 2--3 дня --- высокая скорость
  \item \textbf{Низкая себестоимость:} разработка бота за 2--3 дня, маржа 50--70\%
  \item \textbf{Масштабируемость:} модульная архитектура --- быстрая адаптация под ниши
  \item \textbf{Персонализация:} каждый бот настроен индивидуально (vs. массовые конструкторы)
\end{itemize}
\end{frame}

% === СЛАЙД 6: Технологичность (задел) ===
\begin{frame}{Технологичность стартап-проекта}
\cbar
\textbf{3. Имеющийся задел} (октябрь--ноябрь 2025)

\begin{columns}[T]
\begin{column}{0.48\textwidth}
\textbf{\textcolor{cGray}{Разработка:}}
\begin{itemize}\setlength\itemsep{0pt}
  \item 2 MVP-бота (коучи, товары)
  \item Интеграция с Telegram Bot API
  \item Тестовая техподдержка
  \item Визуальный конструктор ботов
  \item Сайт-лендинг BotWork
\end{itemize}
\end{column}
\begin{column}{0.48\textwidth}
\textbf{\textcolor{cGray}{Рынок:}}
\begin{itemize}\setlength\itemsep{0pt}
  \item 3 Telegram-канала
  \item 5 рекламных постов (8\,000\R)
  \item \textbf{15 лидов, 2 сделки} (13\%)
  \item Анализ 5 конкурентов
  \item Финансовый план
\end{itemize}
\end{column}
\end{columns}
\end{frame}

% === СЛАЙД 7: Технологичность (нацпроект + ИС) ===
\begin{frame}{Технологичность стартап-проекта}
\cbar
\textbf{4. Соответствие направлению нацпроекта}

\medskip
Направление \guillemotleft{}Экономика данных\guillemotright{}:
\begin{itemize}\setlength\itemsep{0pt}
  \item Развитие цифровых сервисов для бизнеса
  \item Автоматизация бизнес-процессов с помощью ИТ
  \item Повышение цифровой зрелости МСП
\end{itemize}

\bigskip
\textbf{5. Интеллектуальная собственность}

\medskip
На текущем этапе зарегистрированная ИС отсутствует.
Планируется регистрация программы для ЭВМ (конструктор ботов) в Роспатенте в первый год проекта.
\end{frame}

% === СЛАЙД 8: Аналоги (существующие) ===
\begin{frame}{Аналоги и конкуренты}
\cbar
\textbf{1. Существующие аналоги}

\begin{itemize}\setlength\itemsep{3pt}
  \item \textbf{BotHelp} --- конструктор автоворонок, 15--50 тыс.\R, сложный интерфейс
  \item \textbf{SendPulse} --- маркетинговая платформа с модулем ботов, от 12 тыс.\R/год, избыточный функционал
  \item \textbf{Salebot} --- конструктор автоворонок, от 10 тыс.\R, для продвинутых пользователей
  \item \textbf{PuzzleBot} --- визуальный конструктор, от 990\R/мес, ограниченный базовый тариф
\end{itemize}

\medskip
\textcolor{cRed}{Все конкуренты предлагают инструменты для самостоятельной сборки. BotWork даёт готовое решение \guillemotleft{}под ключ\guillemotright{} по доступной цене.}
\end{frame}

% === СЛАЙД 9: Аналоги (конкурентные преимущества) ===
\begin{frame}{Аналоги и конкуренты}
\cbar
\textbf{2. Конкурентные преимущества}

\smallskip
\centering
\footnotesize
\begin{tabular}{lcccc}
\toprule
\textbf{Критерий} & \textbf{BotWork} & BotHelp & SendPulse & Salebot \\
\midrule
Цена & \textcolor{cRed}{5--25 тыс.} & 15--50 тыс. & 12 тыс.+/год & 10 тыс.+ \\
Персонализация & \textcolor{cRed}{Под ключ} & Шаблоны & Нет & Частично \\
Поддержка & \textcolor{cRed}{2 нед. бесп.} & Платная & Тикеты & Платная \\
Срок запуска & \textcolor{cRed}{3--7 дней} & Самост. & Самост. & Самост. \\
Конструктор & Да & Да & Да & Да \\
\bottomrule
\end{tabular}

\raggedright
\medskip
\begin{block}{Уникальное ценностное предложение}
Telegram-боты \guillemotleft{}под ключ\guillemotright{} в 2--3 раза дешевле рынка, с бесплатной поддержкой и запуском за 3--7 дней. Готовое решение, а не инструмент для самостоятельного освоения.
\end{block}
\end{frame}

% === СЛАЙД 10: Востребованность ===
\begin{frame}{Аналоги и конкуренты}
\cbar
\textbf{3. Востребованность продукта}

\medskip
\textbf{\textcolor{cGray}{Подтверждённый спрос:}}
\begin{itemize}\setlength\itemsep{0pt}
  \item 15 лидов через 3 Telegram-канала за 2 месяца
  \item 2 сделки (конверсия 13\%), доход 30\,000\R
\end{itemize}

\medskip
\textbf{\textcolor{cGray}{Обратная связь:}}
\begin{itemize}\setlength\itemsep{0pt}
  \item \guillemotleft{}Доступная цена --- не нужно платить 50 тысяч\guillemotright{}
  \item \guillemotleft{}Не нужно разбираться самому --- получил готового бота\guillemotright{}
\end{itemize}

\medskip
\textbf{\textcolor{cGray}{Объективные факторы:}}
\begin{itemize}\setlength\itemsep{0pt}
  \item 50--70 тыс.\ предпринимателей используют Telegram для продаж
  \item Рост интереса к автоматизации: +40\% запросов (2024--2025)
  \item MAU Telegram в РФ более 100 млн
\end{itemize}
\end{frame}

% === СЛАЙД 11: Рынок (область применения) ===
\begin{frame}{Рынок и потребитель}
\cbar
\textbf{1. Область применения продукта}

\begin{itemize}\setlength\itemsep{2pt}
  \item \textbf{Воронки продаж} --- автоматическая цепочка сообщений к покупке
  \item \textbf{Запись на услуги} --- салоны, репетиторы, консультанты
  \item \textbf{Мини-магазины} --- каталог, корзина, онлайн-оплата
  \item \textbf{FAQ-боты} --- автоответы, снижение нагрузки на поддержку
  \item \textbf{Рассылки} --- новинки и акции по базе клиентов
  \item \textbf{Сбор заявок} --- автоприём и маршрутизация менеджерам
  \item \textbf{Чат-поддержка} --- первичная обработка с передачей оператору
\end{itemize}

\smallskip
\textcolor{cGray}{\small Дополнительно: геймификация, лояльность, отзывы, корпоративные боты.}
\end{frame}

% === СЛАЙД 12: Рынок (TAM-SAM-SOM) ===
\begin{frame}{Рынок и потребитель}
\cbar
\textbf{2. Рынок, сегмент, потребитель}

\smallskip
\textbf{TAM:} 4,5 млн МСП в РФ, из них около 500 тыс.\ потенциально заинтересованы.

\textbf{SAM:} 50--70 тыс.\ предпринимателей, активно продающих через Telegram (РФ, СНГ).

\textbf{SOM:} 50--100 клиентов в первый год, 300+ к третьему.

\medskip
\textbf{\textcolor{cGray}{Портрет потребителя:}}
Владелец малого бизнеса (ИП/ООО), 25--45 лет, оборот до 5 млн\R/год. Ниши: образование, услуги, товары. Ищет доступную автоматизацию.

\medskip
\textbf{\textcolor{cGray}{Тренды:}} Telegram-коммерция +40\%/год. MAU Telegram в РФ более 100 млн.

\textbf{\textcolor{cGray}{География:}} Россия, СНГ.
\end{frame}

% === СЛАЙД 13: Рынок (барьеры) ===
\begin{frame}{Рынок и потребитель}
\cbar
\textbf{3. Барьеры для выхода на рынок}

\medskip
\textbf{1. Высокая конкуренция} --- крупные игроки с большими бюджетами.\\
\textcolor{cRed}{Решение: цена в 2--3 раза ниже, персонализация, качественная поддержка.}

\medskip
\textbf{2. Низкий порог входа} --- легко появляются мелкие конкуренты.\\
\textcolor{cRed}{Решение: клиентская база, репутация, повторные продажи (до 30\%).}

\medskip
\textbf{3. Зависимость от Telegram API} --- изменения могут потребовать доработки.\\
\textcolor{cRed}{Решение: мониторинг API, модульная архитектура для быстрой адаптации.}

\medskip
\textbf{4. Ограниченный маркетинговый бюджет.}\\
\textcolor{cRed}{Решение: бесплатные каналы, микро-инфлюенсеры (2--5 тыс.\R\ за пост).}
\end{frame}

% === СЛАЙД 14: План реализации ===
\begin{frame}{План реализации проекта}
\cbar
\begin{columns}[T]
\begin{column}{0.48\textwidth}
\textbf{Направления:}

\smallskip
{\small
\textbf{\textcolor{cGray}{1. Техзадача:}}
MVP $\to$ рассылки/аналитика $\to$ SaaS

\textbf{\textcolor{cGray}{2. Ресурсы:}}
Команда 3 чел., хостинг, домен

\textbf{\textcolor{cGray}{3. Затраты:}}
0м: 20к | 3м: 45к | 12м: 60к

\textbf{\textcolor{cGray}{4. Команда:}}
Старт: 3 | 12м: +1 разработчик

\textbf{\textcolor{cGray}{5. Доход:}}
Разовая оплата 5--25 тыс.\R\ + доработки
}
\end{column}
\begin{column}{0.48\textwidth}
\textbf{Таймлайн:}

\smallskip
{\small
\textbf{\textcolor{cOrange}{Q4 2025 (старт):}}\\
2 MVP, конструктор, сайт, 2 продажи\\
\textcolor{cRed}{Доход 30к | Расходы 20к}

\medskip
\textbf{\textcolor{cOrange}{Q1 2026 (3 мес):}}\\
Рассылки, аналитика, платежи, 8 ботов\\
\textcolor{cRed}{Доход 120к | Расходы 45к}\\
Конверсия 15\%

\medskip
\textbf{\textcolor{cOrange}{Q3 2026 (12 мес):}}\\
SaaS-конструктор, 25 ботов\\
\textcolor{cRed}{Доход 300к | Расходы 60к}\\
Конверсия 25\%, повторные 30\%
}
\end{column}
\end{columns}
\end{frame}

% === СЛАЙД 15: Квалификация (1/2) ===
\begin{frame}{Квалификация заявителя}
\cbar
\textbf{1. Команда проекта}

\smallskip
\textbf{\textcolor{cPurple}{Артемий}} --- тех.\ лидер: разработка ботов, интеграции, техподдержка.\\
{\small Дисциплины: Python, веб-разработка, базы данных.}

\smallskip
\textbf{\textcolor{cPurple}{Денис}} --- маркетинг и продажи: продвижение, лидогенерация, переговоры.\\
{\small Дисциплины: маркетинг, управление проектами.}

\smallskip
\textbf{\textcolor{cPurple}{Борис}} --- старший разработчик: конструктор, боты.\\
{\small Дисциплины: программирование, UX/UI.}

\medskip
\textbf{2. Опыт по теме заявки}
\begin{itemize}\setlength\itemsep{0pt}
  \item 2 MVP-бота для реальных клиентов
  \item Визуальный конструктор чат-ботов
  \item 2 сделки, доход 30\,000\R
  \item Анализ рынка и 5 конкурентов
\end{itemize}
\end{frame}

% === СЛАЙД 16: Квалификация (2/2) ===
\begin{frame}{Квалификация заявителя}
\cbar
\textbf{3. Участие в программе \guillemotleft{}Стартап как диплом\guillemotright{}}

\smallskip
\textcolor{cGray}{\small [Указать майнор, если применимо]}

\bigskip
\textbf{4. Участие в акселераторах, бизнес-инкубаторах}

\smallskip
\textcolor{cGray}{\small [Указать акселератор HSE, если применимо]}

\bigskip
\textbf{5. Участие в ПИШ / молодежной лаборатории / СКБ}

\smallskip
\textcolor{cGray}{\small [Указать, если применимо]}
\end{frame}

% === СЛАЙД 17: Спасибо за внимание ===
\begin{frame}[plain]
\centering
\vspace{1.5cm}
{\Huge\bfseries\textcolor{cPurple}{Спасибо за внимание!}}

\bigskip
\cbar

\bigskip
{\large BotWork --- Telegram-боты для бизнеса}

\bigskip
\begin{tabular}{rl}
  \textbf{Почта:} & [email] \\
  \textbf{Телефон:} & [телефон] \\
  \textbf{Telegram:} & @borysikkk \\
\end{tabular}

\vfill
{\small Номер заявки: \_\_\_\_\_\_\_ \quad ФИО: \_\_\_\_\_\_\_\_\_\_\_\_\_}
\end{frame}

\end{document}
